\par\Needspace{10\baselineskip}\pdfbookmark[1]{Tafel zur bequemeren Berechnung. 1812.}{gauss1812log}
\begin{center}
\rule{0.66\linewidth}{0.3pt}\\
Monatliche Correspondenz, herausgegeben von Freiherrn von \textsc{Zach}. 1812. Nov.\\
\rule{0.66\linewidth}{0.3pt}
\end{center}

\emph{Tafel zur bequemeren Berechnung des Logarithmus der Summe oder Differenz zweier Größen, die selbst nur durch ihre Logarithmen gegeben sind. Von Herrn Professor} \textsc{Gauss}.

Je weiter die Rechenaufgaben der Astronomen beständig zunehmen, desto wichtiger wird für sie jede Erleichterung, auch wenn sie für sich genommen nur klein ist. Die \emph{Monatliche Correspondenz} hat sich darin schon vielfach verdient gemacht, indem sie verschiedene Tafeln aufgenommen hat, deren geringer Umfang eine gesonderte Veröffentlichung nicht erlaubte. Ich veröffentliche darin daher gern eine kleine Tafel, die zwar nicht eigentlich astronomisch ist, aber doch besonders den rechnenden Astronomen willkommen sein wird. Sie könnte künftig sehr zweckmäßig mit einem neuen Druck der kleinen Tafeln von \textsc{La
Lande} verbunden werden. Die Aufgabe, die sie erleichtern soll, tritt bei astronomischen Rechnungen ständig auf; sonst erfordert sie ein dreimaliges oder, wenn man eine einfache Umformung anwendet, doch notwendigerweise ein zweimaliges Nachschlagen in den Logarithmentafeln. Hier wird dies auf ein einziges Nachschlagen reduziert. Die Idee dazu hat meines Wissens \textsc{Leonelli} zuerst angegeben; er wollte eine solche Tafel jedoch für Rechnungen mit 14 Dezimalstellen aufstellen, und gerade dies kann ich nicht für zweckmäßig halten. In diesem Umfang würde sie einen großen Folioband füllen, ihre Berechnung würde ungeheuer viel Arbeit und Zeit erfordern, und sie wäre fast nie und immer nur von geringem Nutzen. Denn so genaue Rechnungen kommen so selten vor---in der eigentlichen praktischen Astronomie nie---, dass die verhältnismäßig doch nur kleine Erleichterung weder die Aufstellung noch auch nur den Kauf einer solchen Tafel lohnen würde. Ich habe diese Tafel schon vor vielen Jahren für meinen eigenen Gebrauch bei Rechnungen mit 5 Dezimalstellen aufgestellt, die in der Praxis am häufigsten sind. Die Erleichterung, die zwar \emph{jedes Mal} klein, bei vielen Tausend Wiederholungen aber sehr beträchtlich ist, hat mir die aufgewandte Mühe bereits reichlich vergolten. Es wäre zu wünschen, dass jemand die Arbeit auf sich nähme, eine ähnliche Tafel in 10 oder 100 Mal so großem Umfang für Rechnungen mit 7 Dezimalstellen zu \SourceLeaf{254}
erstellen, die den gewöhnlichen Logarithmentafeln als sehr wertvolle Ergänzung beigefügt werden könnte.

Die Anordnung der aus drei Spalten bestehenden Tafel ist sehr einfach. Die erste Spalte \(A\) geht von 0 bis 2 in Schritten von je einem Tausendstel, von dort bis \(3,4\) in Hundertsteln und von \(3,4\) bis \(5,0\) in Zehnteln. Mit \(5,0\) kann die Tafel für 5 Dezimalstellen als abgeschlossen gelten, da die zweite Spalte für diesen und für größere Werte von \(A\) verschwindet und die Zahlen der dritten Spalte denen der ersten gleich werden. Schreibt man eine Zahl der ersten Spalte als \(A=\log m\), so steht in der zweiten Spalte \(B=\log(1+\frac{1}{m})\) und in der dritten Spalte \(C=\log(1+m)\), sodass stets \(C=A+B\) gilt. Man kann die Zahlen der drei Spalten daher auch als das Doppelte der Logarithmen der Tangenten, Kosekanten und Sekanten der Winkel von \(45^{0}\) bis \(90^{0}\) betrachten. Die Anwendung ist folgende:

\begin{center}
I.\quad \emph{Aus den Logarithmen zweier Größen \(a,b\) den Logarithmus ihrer Summe bestimmen.}
\end{center}

Es sei \(\log a\) der größere Logarithmus. Mit \(\log a-\log b\) geht man in die Spalte \(A\) und entnimmt daneben entweder der zweiten Spalte \(B\) oder der dritten Spalte \(C\). Dann gilt
\[
\begin{aligned}
&\log(a+b)=\log a+B\\
\text{oder}\qquad &\log(a+b)=\log b+C.
\end{aligned}
\]

\Needspace{13\baselineskip}\begin{center}
II.\quad \emph{Aus den Logarithmen zweier Größen \(a,b\) den Logarithmus ihrer Differenz bestimmen.}
\end{center}

\emph{Erstens:} Ist die Differenz der Logarithmen \(\log a-\log b\) größer als \(0,30103\), so sucht man sie in \(C\). Damit erhält man
\[
\begin{aligned}
&\log(a-b)=\log a-B\\
\text{oder}\qquad &\log(a-b)=\log b+A
\end{aligned}
\]

\emph{Zweitens:} Ist \(\log a-\log b\) kleiner als \(0,30103\), so sucht man dies in \(B\). Damit erhält man
\[
\begin{aligned}
&\log(a-b)=\log a-C\\
\text{oder}\qquad &\log(a-b)=\log b-A
\end{aligned}
\]
\SourceLeaf{255}
Für jede Aufgabe gibt es daher zwei Lösungswege. Es ist jedoch ratsam, sich an einen bestimmten zu gewöhnen, damit die Benutzung der Tafel leichter zur Gewohnheit wird. Für mich war dabei die jeweils zuerst angegebene Methode am bequemsten.

\Needspace{12\baselineskip}\begin{center}\textit{Beispiele:}\end{center}

I. Um aus $\log a=0{,}36173$ und $\log b=0{,}23045$ den Logarithmus der Summe zu bestimmen, sucht man $0{,}13128$ in $A$. Dabei erhält man
\[
\begin{array}{rcr}
B&\cdot\;\cdot\;\cdot\;\cdot\;\cdot&0{,}24033\\
\log a&\cdot\;\cdot\;\cdot\;\cdot\;\cdot&0{,}36173\\\hline
\log(a+b)&\cdot\;\cdot&0{,}60206
\end{array}
\qquad
\begin{array}{rcr}
C&\cdot\;\cdot\;\cdot\;\cdot\;\cdot&0{,}37161\\
\log b&\cdot\;\cdot\;\cdot\;\cdot\;\cdot&0{,}23045\\\hline
&&0{,}60206
\end{array}
\]

II. Aus $\log a=0{,}89042$ und $\log b=0{,}24797$ soll der Logarithmus der Differenz bestimmt werden. Da $\log a-\log b=0{,}64245$ größer als $0{,}30103$ ist, sucht man dies in der Spalte $C$; daneben findet man
\[
\begin{array}{rcr}
B&\cdot\;\cdot\;\cdot\;\cdot\;\cdot&0{,}11227\\
\log a&\cdot\;\cdot\;\cdot\;\cdot\;\cdot&0{,}89042\\\hline
\log(a-b)&\cdot\;\cdot&0{,}77815
\end{array}
\qquad
\begin{array}{rcr}
A&\cdot\;\cdot\;\cdot\;\cdot\;\cdot&0{,}53018\\
\log b&\cdot\;\cdot\;\cdot\;\cdot\;\cdot&0{,}24797\\\hline
&&0{,}77815
\end{array}
\]

III. Aus $\log a=0{,}25042$ und $\log b=0{,}19033$ soll der Logarithmus der Differenz bestimmt werden. Hier ergibt $\log a-\log b=0{,}06009$, in $B$ aufgesucht,
\[
\begin{array}{rcr}
C&\cdot\;\cdot\;\cdot\;\cdot\;\cdot&0{,}88871\\
\log a&\cdot\;\cdot\;\cdot\;\cdot\;\cdot&0{,}25042\\\hline
\log(a-b)&\cdot\;\cdot&9{,}36171
\end{array}
\qquad
\begin{array}{rcr}
A&\cdot\;\cdot\;\cdot\;\cdot\;\cdot&0{,}82862\\
\log b&\cdot\;\cdot\;\cdot\;\cdot\;\cdot&0{,}19033\\\hline
&&9{,}36171
\end{array}
\]

